\section{Reflexion}

Ein besonderes Merkmal dieses Themas ist seine visuelle Anziehungskraft. Die Untersuchung von Projektionen hochsymmetrischer Körper bietet eine unmittelbare, ästhetische Zugänglichkeit. Obwohl die zugrundeliegende Mathematik – mit Symmetriegruppen, Matrixalgebra und Orthogonalprojektionen – ein sehr hohes und abstraktes Anforderungsniveau besitzt, verhindert die visuelle Aufarbeitung, dass die Theorie trocken wirkt, und bietet einen sehr konkreten Zugang zur Materie. Diese Dualität macht es zu einem äußerst dankbaren Thema: Die Lernenden oder Lesenden können die abstrakten algebraischen Beweise durch die geometrischen Projektionsbilder (wie die Umrisse des Oktaeders) stets kontrollieren und verifizieren. Hier wäre wahrscheinlich auch das grösste Potential vorhanden gewesen, um digitale Medien gezielter einzusetzen. Dieses Thema bietet sich sehr dazu an visuell zu arbeiten und verschiedene Medien, wie Grafiken, Applets, Videos, einzusetzen, um die Dinge zu visualisieren.
\newline
\newline
Zwar ist das vorliegende Skript bewusst in einem sehr formalen, mathematisch rigorosen Stil verfasst und dementsprechend anspruchsvoll, nichtsdestotrotz schließt dies eine Behandlung im gymnasialen Unterricht keinesfalls aus. Im Gegenteil: Gerade durch die hohe visuelle Zugänglichkeit lässt sich die komplexe Theorie gut auf ein greifbares Niveau herunterbrechen. Während für das vollständige Durchdringen der algebraischen Beweise eine gewisse mathematische Reife erforderlich ist, sind die geometrischen Kernaussagen intuitiv erfassbar und können auch auf gymnasialer Stufe problemlos verstanden werden.  
\newline
\newline
Das Thema bietet zudem die seltene Gelegenheit, tiefgreifende lineare Algebra mit einem handlungsorientierten Ansatz zu verknüpfen. Die Möglichkeit, die untersuchten Polyeder (wie die platonischen Körper oder das (Pseudo-)Rhombenkuboktaeder) selbst aus Papier oder Modellbausätzen zu basteln, schlägt eine Brücke zur haptischen Erfahrung. Somit könnte dieses Thema auch Schülerinnen und Schüler (SuS) ansprechen, die ansonsten keinen großen Bezug zur klassischen, rein kalkülbasierten Schulmathematik haben. 
\newline
\newline
Im klassischen Mathematikstudium werden Lineare Algebra und Geometrie häufig sehr isoliert voneinander gelehrt. Lineare Algebra verbleibt oft in abstrakten Vektorräumen, während sich dieses Thema genau an der Schnittmenge bewegt. Es zeigt eindrucksvoll auf, welche enorme geometrische Aussagekraft algebraische Werkzeuge (Symmetrische Matrizen, Diagonalisierbarkeit, Eigenwerte) haben. Ein weiterer didaktisch wertvoller Aspekt der Arbeit ist der starke Bezug zur Darstellenden Geometrie, einem Fachbereich, der in den aktuellen Schul-Curricula oftmals in den Hintergrund geraten ist. Wenn auch anspruchsvoll, kann dieses Thema als eine Weiterführung in diesem Zusammenhang behandelt oder erwähnt werden.
\newline
\newline
Zuletzt bietet das Thema ein reiches Potenzial für die Integration der Mathematikgeschichte. Die Faszination für hochsymmetrische Körper reicht von den antiken Philosophen (Platonische Körper) über Archimedes bis hin zu Johannes Kepler, der diese Formen in seinem kosmologischen Modell des Sonnensystems nutzte. Gleichzeitig spannt das Pseudo-Rhombenkuboktaeder den Bogen in die Moderne, da es erst im 20. Jahrhundert im Rahmen der Johnson-Körper systematisch Beachtung fand. Diese historische Dimension verleiht der Mathematik ein menschliches, kulturelles Gesicht und zeigt, dass die Geometrie ein stetig wachsendes, lebendiges Forschungsfeld ist.